% ============================================================
% Chapter 6 — The Four Learning Interventions
% ============================================================
\chapter{The Four Learning Interventions}
\label{ch:interventions}

\roadmap{This chapter treats each of the four learning-strategy
interventions in the same order---theory, mechanics, engineering---so that
its pedagogical warrant, its learner-facing behavior, and its
implementation are visible together. It then documents three
cross-cutting mechanisms: how a text selection reaches the generator, how
the chatbot suggests strategies, and how teachers configure prompts. It
closes by establishing, against the code, that interventions are
exclusively learner-initiated and never triggered by affect.}

\section{Common structure}
\label{sec:interventions-common}

All four interventions share a lifecycle. A learner initiates one---by
clicking a strategy button or by accepting a chatbot suggestion---
optionally after selecting text on the page. The backend resolves the
grounding context (\cref{sec:modes-context}), calls the LLM with a
system prompt (teacher-configurable) and a hard-coded user-prompt
template, and stores the generated activity as a
\dbtable{LearningIntervention} row whose \code{sessionData} JSON holds the
activity state.\footnote{\code{apps/api/src/learning-interventions};
\dbtable{LearningIntervention}, \dbtable{InterventionType}} Interaction
then proceeds through intervention-specific endpoints. A learner may save
any activity for later review (\dbtable{SavedInterventionReview}), which
surfaces in the review queue. \Cref{tab:interventions-summary} orients the
four before the detailed treatment.

\begin{table}[htbp]
\centering
\small
\caption{The four interventions at a glance.}
\label{tab:interventions-summary}
\begin{tabular}{@{}p{0.19\linewidth}p{0.25\linewidth}p{0.22\linewidth}p{0.24\linewidth}@{}}
\toprule
Intervention & Generates & Grading & Persistence \\
\midrule
Practice testing & Mixed MCQ + short answer (default 3+2) & MCQ
deterministic; short answer LLM with keyword fallback & \dbtable{LearningIntervention} \\
Interrogative elaboration & ``Why/how'' question suggestions + dialogue &
None (LLM depth summary) & \dbtable{LearningIntervention} \\
Stepwise learning & 3--6 scaffolded steps with checks & LLM per step; force-pass after 2 attempts & \dbtable{LearningIntervention} \\
Distributed practice & Flashcards on an SM-2 schedule & Learner self-rating drives SM-2 & \dbtable{LearningIntervention} + \dbtable{SpacedRepetitionCard} \\
\bottomrule
\end{tabular}
\end{table}

\section{Practice testing}
\label{sec:interventions-pt}

\paragraph{Theory.} Retrieval practice is a high-utility technique
\citep{dunlosky2013}: the act of retrieving strengthens later retention
more than restudy, even when it feels less fluent
\citep{roediger2006, bjork1994}.

\paragraph{Mechanics.} The intervention generates a mixed set of
multiple-choice and short-answer items over the current material
(default three MCQ and two short answer), which the learner answers and
submits for immediate scoring.

\paragraph{Engineering.} Generation and submission are served by dedicated
endpoints (\cref{app:endpoints}); the item set and, after submission, the
answers and score live in \code{sessionData}. Grading is
mixed.\footnote{\code{learning-interventions.service.ts}: \code{gradeAnswer}
(deterministic), \code{gradeShortAnswerWithLlm} (LLM).} MCQ answers are
normalized to a single letter and compared deterministically; non-empty
short answers are graded by an LLM under a hard-coded ``be generous''
rubric, falling back to keyword matching (at least half the expected
keywords) if the LLM call fails. The final score is the percentage
correct. One consequence worth stating: because short-answer grading is
LLM-based, identical answers may occasionally receive different scores
across runs.

\section{Interrogative elaboration}
\label{sec:interventions-ie}

\paragraph{Theory.} Elaborative interrogation---prompting learners to
answer \emph{why} and \emph{how} rather than \emph{what}---has
moderate-utility evidence \citep{dunlosky2013}, and self-generated
elaboration aids memory through the generation effect
\citep{pressley1987}.

\paragraph{Mechanics.} The intervention suggests a set of why/how
questions grounded in the material and then hosts an explanatory
dialogue: the tutor's role is to explain, not to quiz. On completion it
produces a descriptive depth rating (surface/moderate/deep), not a score.

\paragraph{Engineering.} Suggestion, per-turn asking, and completion are
separate endpoints; the conversation and its per-turn selection context
accumulate in \code{sessionData}. There is no correctness grading. A
historical bug---context frozen at the first slide---was fixed by
re-resolving grounding on each turn, and per-turn selection and page are
now persisted so the review surface can show what each question referred
to.\footnote{\code{docs/BUG_REPORT_20260612.md}}

\section{Stepwise learning}
\label{sec:interventions-sl}

\paragraph{Theory.} Stepwise learning operationalizes scaffolding
\citep{wood1976, vandepol2010} within the zone of proximal development
\citep{vygotsky1978}, decomposing material to bound cognitive load
\citep{sweller1988} and fading support as competence is shown.

\paragraph{Mechanics.} The intervention breaks the material into three to
six progressive steps, each with a comprehension check; the learner
advances step by step.

\paragraph{Engineering.} Generation, per-step checking, and advancement
are separate endpoints; per-step attempts and outcomes are held in
\code{sessionData}. Checking is LLM-based with no deterministic fallback.
A design choice is disclosed for the record: after two failed attempts on
a step the correct answer is shown and the step is force-marked passed, so
a learner is never trapped---at the cost that ``passed'' does not imply
``answered correctly.''

\section{Distributed practice}
\label{sec:interventions-dp}

\paragraph{Theory.} Distributed (spaced) practice is the second
high-utility technique \citep{dunlosky2013, cepeda2006}; scheduling
reviews at expanding intervals exploits the spacing effect
\citep{bjork1994}.

\paragraph{Mechanics.} The intervention generates flashcards, which the
learner reviews and self-rates; ratings drive an SM-2 schedule
\citep{wozniak1994} that sets each card's next due date.

\paragraph{Engineering.} Each card is a \dbtable{SpacedRepetitionCard} row
seeded with ease $2.5$, interval $0$, and repetitions $0$; reviews update
these fields. The scheduler is the SM-2 algorithm, reproduced verbatim
from the implementation in \cref{lst:sm2}.

\begin{lstlisting}[caption={SM-2 update, verbatim from
\code{utils/sm2-algorithm.ts}. Rating map: again=1, hard=2, good=3,
easy=5.}, label={lst:sm2}]
if (quality >= 3) {
  if (repetitions === 0) interval = 1;
  else if (repetitions === 1) interval = 6;
  else interval = Math.round(interval * ease);
  repetitions += 1;
} else {
  repetitions = 0;
  interval = 1;
}
ease = Math.max(1.3, ease + (0.1 - (5 - quality) * (0.08 + (5 - quality) * 0.02)));
// nextReviewAt = today + interval days
\end{lstlisting}

The exact constants matter for reproducibility and are stated in prose:
the ease factor is floored at $1.3$ and seeded at $2.5$; intervals expand
$1 \rightarrow 6 \rightarrow \operatorname{round}(\text{interval} \times
\text{ease})$; a rating below ``good'' (quality $<3$) is a full lapse,
resetting repetitions to $0$ and interval to $1$. Two implementation
notes distinguish this from textbook SM-2: the rating map never emits
quality $4$ (``easy'' jumps to $5$), and ``hard'' (quality $2$) is treated
as a full lapse rather than a shortened interval. Both are faithful to
the code and are flagged so that any recall-calibration analysis
(\cref{sec:methodology-analyses}) interprets the schedule correctly.

\section{Text selection to generator}
\label{sec:interventions-selection}

A learner grounds an intervention on a passage by selecting it. Selection
is captured on two surfaces: on block pages, a mouse-up selection is
accepted only if it exceeds 20 characters and originates inside a
selectable element; in the PDF reader, a selection above 10 characters is
captured with its page
number.\footnote{\code{apps/web/src/components/.../PageContext.tsx};
\code{PdfReader.tsx}} The selection travels in the request body and is
resolved server-side by the 20-character rule of
\cref{sec:modes-context}, becoming the grounding text substituted into the
prompt template. The authoritative threshold is therefore 20 characters;
the PDF surface's 10-character gate governs capture only.

\section{The strategy-suggestion engine}
\label{sec:interventions-suggest}

The docked chatbot may suggest a strategy, but the ``engine'' is a
\emph{prompt instruction}, not a rule
system.\footnote{\code{learning-interventions.service.ts:4214}} The
persona instructs the model to append, on its final line, a tag
\code{[SUGGEST:KEY]} with \code{KEY} one of the four intervention types,
only when genuinely relevant. A deterministic regular expression parses
the tag, strips it from the visible reply, and stores the parsed value in
\dbtable{ChatbotMessage.suggestedStrategy}. The suggestion renders as a
clickable card; the intervention launches only when the learner clicks it.
Persisting the suggested strategy lets later analysis measure suggestion
uptake against actual intervention initiation.

\section{Teacher prompt configuration}
\label{sec:interventions-teacher}

Teachers customize interventions through \dbtable{InterventionPromptConfig},
which stores a per-course \emph{system prompt} per intervention type and,
for practice testing only, default MCQ and short-answer counts.
Editability is deliberately bounded: the system prompt is teacher-editable
(validated for length and warned if it omits JSON-format instructions),
but the \emph{user-prompt template}, all grading and summary prompts, the
chatbot persona, and structural constants (count clamps, the two-attempt
cap, the SM-2 constants) are hard-coded. The split---editable system
prompt, fixed templates---keeps teacher customization from breaking the
machine-readable output contracts on which grading depends.

\section{Interventions are learner-initiated, never affect-triggered}
\label{sec:interventions-autotrigger}

A claim central to this report's methodology is that the sensing layer
does not drive the pedagogy. The codebase supports it without
qualification. Interventions carry only two trigger reasons anywhere in
the service, \code{student_initiated} and \code{pre_generated}; no
affective, biometric, or automatic trigger reason exists. The
\code{INTERVENTION_TRIGGERED} activity action is a logging event written
\emph{after} a learner's request, not a trigger. Every generation
endpoint is reachable only from a learner action---a strategy button or an
accepted suggestion card. The affect and text-mining pipelines are
consumers, not producers, of these events: the only cross-link runs the
other way, with the chatbot feeding learner utterances into text mining.
The closed loop---affect-triggered support, the natural endgame of
RQ2---is therefore \emph{not} implemented, and is reserved for future
work (\cref{ch:future}). Keeping the loop open is what allows RQ1's
detection validity to be established observationally, free of the
confound that would arise if the intervention were triggered by the very
state being measured.

\medskip
\noindent The pedagogy layer complete, the next chapter documents the
sensing and logging layer that observes learners as they use it.
